\documentclass[main]{subfiles}


\title[Dynamic Approaches]{Dynamic Approaches}
\subtitle{DAC by RL}
\author[André Biedenkapp]{André Biedenkapp}
\date{}

\titlebackground*{images/AutoML_LearningToLearn2}

%\institute{}
%\date{}

\begin{document}
	
	\maketitle
%-----------------------------------------------------------------------
%-----------------------------------------------------------------------
%----------------------------------------------------------------------
\begin{frame}{Dynamic Algorithm Configuration by RL}
    \centering
    \begin{tikzpicture}[node distance=4cm]
        %PreProcessing
        
        \node (Agent) [activity,
        minimum height=.9cm] {Policy $\policy$};
        
        \node (Algo) [activity, right of=Agent, xshift=3cm,
        minimum height=.9cm] {Algorithm};
        
        \begin{pgfonlayer}{background}
        \path (Agent -| Agent.west)+(-0.12,1.75) node (resUL) {};
        \path (Algo.east |- Algo.south)+(0.125,-1.75) node (resBR) {};
        
        % Context
        \onslide<2->{
        \path [rounded corners, draw=black!50, fill=white] ($(resUL)+(0.5, -0.5)$) rectangle ($(resBR)+(0.5, -0.5)$);
        \path [rounded corners, draw=black!50, fill=white] ($(resUL)+(0.375, -0.375)$) rectangle ($(resBR)+(0.375, -0.375)$);
        \path [rounded corners, draw=black!50, fill=white] ($(resUL)+(0.25, -0.25)$) rectangle ($(resBR)+(0.25, -0.25)$);
        \path [rounded corners, draw=black!50, fill=white] ($(resUL)+(0.125, -0.125)$) rectangle ($(resBR)+(0.125, -0.125)$);
        }
        % Top level
        \path [rounded corners, draw=black!50, fill=white] (resUL) rectangle (resBR);
        \path (resBR)+(-.9,0.175) node [text=black!75] {dataset $\dataset$};
        \end{pgfonlayer}
        
        %        \draw[myarrow] (feat.south) -- ($(feat.south |- Agent)+(0,1.125)$);
        
        \draw[myarrow] (Agent.north) -- ($(Agent.north)+(0.0,+0.5)$) -- ($(Algo.north)+(0.0,+0.5)$) node [above=-.05,pos=0.5] { action $\policy(s_t) \to a_t$} node [below,pos=0.5] {(set param $\conf$)} -- (Algo.north);
        \draw[myarrow, dashed] ($(Algo.south)+(-0.25, 0)$) -- ($(Algo.south)+(-0.25, -0.5)$) -- ($(Agent.south)+(0.25, -0.5)$) node [above,pos=0.5] {state $s_{t+1}$} -- ($(Agent.south)+(0.25, 0)$);
        \draw[myarrow] ($(Algo.south)+(0.25, 0)$) -- ($(Algo.south)+(0.25, -0.95)$) -- ($(Agent.south)+(-0.25, -0.95)$) node [below,pos=0.5] {reward $r_{t+1}$} -- ($(Agent.south)+(-0.25, 0)$);
        
        \onslide<2->{
        \draw[<->, thick, draw=black!32.5] (resBR.east |- resUL.center) -- ($(resBR.east |- resUL.center)+(0.5, -0.5)$);
        \path (resBR.east |- resUL.center)+(.5, -0.175) node [text=black!32.5] {$\datasets$};
        % This path is only needed so the figure stays centered (the text is not visible)
        \path (resUL.west |- resUL.center)+(-.5, -0.175) node [text=black!0] {$\datasets$};}
        
    \end{tikzpicture}
    \begin{itemize}
        \item RL is not setup with generalization across problems (i.e. environments) in mind!
        \item<2-> To be able to leverage the contextual MDP we have to slightly modify the training and policy
    \end{itemize}
\end{frame}
%-----------------------------------------------------------------------
%----------------------------------------------------------------------
\begin{frame}{General Policies}
    \begin{itemize}
        \item We want to learn a policy that performs well on different datasets (i.e., the context)\pause
        \item We can choose two different approaches \liturl{Biedenkapp 2026}{https://andrebiedenkapp.github.io/assets/pdf/paper/26-AAMAS-Blue-Sky.pdf}:\pause
        \begin{itemize}
            \item Context oblivious policies $\policy(s)$ purely observe the state of the algorithm that is being tuned
            \begin{itemize}\pause
                \item Context oblivious policies generalize by training on a distribution of datasets.
                \item Context oblivious policies are robust and tolerate variation but might can not perfectly adapt to a specific dataset.
                \item Simpler to realize
            \end{itemize}\pause
            \item Context aware policies $\policy(s,\dataset)$ that additionally observe features of the dataset at hand
            \begin{itemize}\pause
                \item Context aware policies can directly adapt to the dataset at hand
                \item Require additional design decisions on A) what is an appropriate context and B) how to best include it in the learning process.
            \end{itemize}
        \end{itemize}
    \end{itemize}
\end{frame}
%-----------------------------------------------------------------------
%----------------------------------------------------------------------
\begin{frame}{Actions}
Domains of change
    \begin{itemize}
        \item Directly setting the hyperparameter $\conf_{i,t} = \action_t$
        \begin{itemize}\pause
            \item Full flexibility for adaptation
            \item Can lead to very bumpy schedules (i.e. bang-bang control)
        \end{itemize}\pause
        \item Relative updates
        \begin{itemize}\pause
            \item Change the hyperparameter by a factor, e.g., $\conf_{i,t} = \action_t\cdot\conf_{i,t-1}$
            \item Change the hyperparameter by a summand, e.g., $\conf_{i,t} = \action_t+\conf_{i,t-1}$
            \item Less flexible but likely less bumpy schedules
        \end{itemize}
    \end{itemize}
\end{frame}
%-----------------------------------------------------------------------
%----------------------------------------------------------------------
\begin{frame}{States}
    \begin{itemize}
        \item Start with everything that is used in existing manually designed heuristics\pause
        \item For example, for learning rates:
        \begin{itemize}\pause
            \item Previous learning rate(s)
            \item Number of updates
            \item Loss
            \item Momentum
            \item Gradient information
        \end{itemize}\pause
        \item Why are these good features?
        \begin{itemize}\pause
            \item already informative to update manually designed policies
            \item already available in existing implementations
            \item cheap to compute
            \item (indirectly) quantify the progress of the target algorithm
        \end{itemize}\pause
        \item If these are very noisy statistics, consider to smooth them over the last updates
    \end{itemize}
\end{frame}
%-----------------------------------------------------------------------
%----------------------------------------------------------------------
\begin{frame}{Transitions}
    \begin{itemize}
        \item Implicitly given by training the algorithm/model for an iteration/a few steps\pause
        \item The Granularity here can be very important\footnote{PBT-style training also requires careful consideration of this detail}
        \begin{itemize}\pause
            \item Too often: Can lead to extremely difficult learning problems
            \item Too infrequent: We might potentially miss some important updates and not update the hyperparameter(s) accordingly
        \end{itemize}\pause
        \item If you assume that your hyperparameters can stay constant for quite some time (without
knowing exactly for how many updates), you can try to use an RL approach that is able
to efficiently repeat actions, e.g. \liturl{Biedenkapp et al. 2021}{https://proceedings.mlr.press/v139/biedenkapp21a.html}.
    \end{itemize}
\end{frame}
%-----------------------------------------------------------------------
%----------------------------------------------------------------------
\begin{frame}{Reward}
    \begin{itemize}
        \item We minimize cost (e.g., validation loss), but RL agents maximize reward $r(s)=-\cost(s)$
        \item RL agents optimize (the expected discounted) cumulative reward
        \begin{itemize}
            \item We go quickly for better cost values during training
            \item Often fairly well correlated with final loss
            \item BUT: Can lead to too aggressive hyperparameter changes early on to minimize cost quickly, but does not convergence as well later on
            \item You can play with the discount factor to account for this
        \end{itemize}
    \end{itemize}
\end{frame}
%-----------------------------------------------------------------------
%----------------------------------------------------------------------
\begin{frame}{Context-Awareness: Characterizing Datasets}
    \begin{itemize}
        \item An RL agent could try to implicitly learn what kind of dataset $\dataset$ is at hand right now\pause
        \begin{itemize}
            \item Possible, given a sufficiently large history of observations
            \item Might require a lot of training and large history
        \end{itemize}\pause
        \item Directly provide meta-features characterizing the dataset
        \begin{itemize}
            \item Counting features, such as, number of observations, number of features, number of classes
            \item Landmarking features, such as, loss of tree-based models vs. linear model
        \end{itemize}\pause
        \item Simplest approach to have a policy context-aware: Concatenation to the state-observations \liturl{Hallak et al. 2015}{https://arxiv.org/abs/1502.02259}
        \item This simple approach often does always outperform context-oblivious policies! See, e.g, \liturl{Benjamins et al. 2023}{https://openreview.net/forum?id=Y42xVBQusn} or \liturl{Prasanna et al. 2024}{https://openreview.net/forum?id=o8DrRuBsQb}\pause
        \item The choice of appropriate context might itself be problem specific
    \end{itemize}
\end{frame}
%-----------------------------------------------------------------------
%----------------------------------------------------------------------
\begin{frame}{Main Takeaways} 

%Most important insights from this video:

\begin{itemize}
    \item \textbf{Context-oblivious policies} can learn robust and generalizable behavior by learning across a distribution of environment contexts.
    \item \textbf{Context-aware policies} can directly adapt to the dataset at hand, but require careful consideration on what are appropriate context features.
    \item \textbf{DAC by RL} is one of the most powerful approaches for dynamic tuning but requires potentially much more compute.
\end{itemize}

\end{frame}

\end{document}
